% Generated draft: requires mathematical and source review.
\documentclass[11pt]{article}
\usepackage[margin=25mm]{geometry}
\usepackage{fontspec}
\setmainfont{Latin Modern Roman}
\usepackage{amsmath,amssymb,mathtools,unicode-math}
\setmathfont{Latin Modern Math}
\usepackage{enumitem,needspace,xurl,hyperref}
\setlength{\parindent}{0pt}
\setlength{\parskip}{5pt}
\setlength{\emergencystretch}{3em}
\newcommand{\N}{\mathbb{N}}
\newcommand{\Z}{\mathbb{Z}}
\newcommand{\Q}{\mathbb{Q}}
\newcommand{\R}{\mathbb{R}}
\newcommand{\C}{\mathbb{C}}
\newcommand{\F}{\mathbb{F}}
\newcommand{\A}{\mathbb{A}}
\newcommand{\PP}{\mathbb{P}}
\newcommand{\E}{\mathbb{E}}
\newcommand{\eps}{\varepsilon}
\newcommand{\dd}{\,\mathrm{d}}
\begin{document}
\section*{UnsolvedMath ID: 30005208}
\textbf{Source problem code:} catalogue entry
\subsection{Definitions and mathematical statement}
Let \(X\) denote the ambient manifold intended by the source, and let \(\omega\) and \(g\) denote its symplectic form and Riemannian metric. Since the evolving Lagrangians are surfaces, \(X\) has real dimension four. The ambient geometry and the hypotheses ensuring preservation of the Lagrangian condition are not specified in the supplied excerpt. A smooth Lagrangian mean curvature flow is a family of immersions \(F_t\colon L\to X\) satisfying
\[
F_t^*\omega=0,\qquad (\partial_tF_t)^\perp=H_t,
\]
where \(H_t\) is the mean curvature vector and \(\perp\) denotes normal projection with respect to \(g\).

At a singular space-time point \((x,T)\), a tangent flow means a local weak limit of parabolic rescalings about \((x,T)\). In local coordinates centered at \(x\), these rescalings have the form
\[
L_s^{(j)}=\lambda_j\bigl(L_{T+\lambda_j^{-2}s}-x\bigr),\qquad \lambda_j\longrightarrow\infty,\quad s<0,
\]
where \(L_t=F_t(L)\). This expression is interpreted locally in the tangent space \(T_xX\); subtraction denotes the chosen local coordinates. The weak convergence framework and any continuation through earlier singularities remain unspecified by the record.

The supplied optimistic expectation is the following schematic assertion: for generic admissible initial Lagrangian surfaces, every tangent flow at every singularity is a static union of two distinct multiplicity-one Lagrangian planes. More explicitly, if \((\mu_s)_{s<0}\) is any such tangent flow, there should exist distinct two-dimensional Lagrangian linear subspaces \(P_1,P_2\subset T_xX\) such that
\[
\mu_s=\mathcal H^2\!\llcorner P_1+\mathcal H^2\!\llcorner P_2
\quad\text{for every }s<0,
\qquad
\dim(P_1\cap P_2)\in\{0,1\}.
\]
Here \(\mathcal H^2\!\llcorner P_i\) is two-dimensional Hausdorff measure restricted to \(P_i\), so each plane occurs with multiplicity one. The case \(\dim(P_1\cap P_2)=0\) means that the planes meet only at the rescaling origin, and they are transverse in the four-dimensional tangent space. The case \(\dim(P_1\cap P_2)=1\) means that they meet along a line. This is a formulation of the proposed classification, not an assertion that it holds. In particular, the admissible initial-data class, its topology, and the meaning of ``generic'' must be supplied before this becomes a precise conjecture.

The surrounding excerpt gives additional, incomplete context. It attributes to Thomas--Yau's Conjecture 7.3, cited there as [14], a prediction that a finite-time singularity permits a decomposition \(L=M_1\#M_2\) as a graded connected sum satisfying conditions labeled (1), which are absent from the record. It reports a result supporting one such scenario when the tangent flow at the first singular time consists of two transverse planes, and says that the graded connected-sum decomposition is guaranteed when \(L\) is a sphere in that setting. The full hypotheses of that result are not supplied. Joyce's conjectural picture, cited as [3], allows other singularities, notably tangent flows consisting of two static planes meeting along a line; the excerpt attributes some progress on these singularities to the authors' previous work [5].

The excerpt also mentions a holomorphic volume form \(\Omega\) on \(X\), a Lagrangian angle \(\theta\), and ``cohomological'' Lagrangian angles \(\phi(M_i)\). Under an additional almost-calibrated hypothesis it displays, in incomplete surrounding text,
\[
\phi(M_1),\ \phi(M_2)\subset\bigl(\inf_L\theta,\sup_L\theta\bigr).
\]
This is retained as a transcription of the supplied qualification, not as a fully specified mathematical assertion: the definitions, grading choices, full hypotheses, and whether the quantities \(\phi(M_i)\) are numbers or sets are missing. These omissions prevent recovery of a precise connected-sum or angle theorem from this record.
\subsection{Short English statement}
The proposed generic-singularity picture predicts that every tangent flow of a Lagrangian surface consists of two static planes, each with multiplicity one, meeting either at one point or along a line. The record leaves the ambient hypotheses and the meaning of ``generic'' unspecified, so this remains a schematic conjecture requiring mathematical review. Its connected-sum and angle statements are incomplete contextual qualifications. The supplied status is ``partially solved''; this label and the accompanying literature assessment are dated curation, not independently verified conclusions.
\subsection{Sources}
\begin{itemize}
\item Supplied local UnsolvedMath registry record: permanent ID \(30005208\), exact source problem code \texttt{catalogue entry}, title ``Thomas Yau Singularity Picture for Lagrangian Mean Curvature Flow,'' in the Oberwolfach Reports Open Problems collection. No named proposer is supplied; the proposed year is recorded as 2022.
\item Supplied report citation: \emph{Calculus of Variations (2023)}, Oberwolfach Reports, pp. 2129--2193, DOI 10.4171/owr/2022/37. \url{https://doi.org/10.4171/owr/2022/37}. This is the supplied report-level reference; the exact contribution, authors, and passage location are not identified. The citation's 2023 designation is retained alongside the record's 2022 proposed year and DOI component.
\item Supplied literature source: ``Lotay--Schulze--Szekelyhidi, neck pinches.'' \url{https://arxiv.org/abs/2208.11054}.
\item Supplied literature source: ``Szekelyhidi, generic Lagrangian neck pinches.'' \url{https://arxiv.org/abs/2602.15771}.
\item Supplied literature source: ``original Oberwolfach report.'' \url{https://doi.org/10.4171/owr/2022/37}.
\item The literature check date is exactly \texttt{2026-08-22}. The supplied dated curation labels the problem ``partially solved'' and assesses that the two-plane neck-pinch scenario is understood, that recent work treats circle-symmetric and nondegenerate neck pinches, and that the broader Thomas--Yau singularity picture and generic classification remain unresolved. This assessment is supplied dated curation, not independently verified; its time-sensitive status and cited claims require independent checking.
\item Attribution within the excerpt: Thomas--Yau, Conjecture 7.3 in its reference [14]; Joyce's conjectural picture in its reference [3]; and the unidentified excerpt authors' previous work in its reference [5]. The corresponding bibliography entries are absent and have not been reconstructed.
\end{itemize}
\end{document}
